\chapter{نتیجه‌گیری و کارهای آینده}
\label{chap:conclusion}

\section{نتیجه‌گیری}
در این پایان‌نامه یک چارچوب تست سرتاسری متن‌باز برای نرم‌افزارهای مبتنی بر
\lr{Qt C++}
در بستر ویندوز طراحی و پیاده‌سازی شد. چارچوب روی زیرساخت استاندارد
\lr{WinAppDriver}/\lr{Appium}
سوار است و آنچه این زیرساخت نمی‌دهد — پایداری، خوانایی و نگهداشت‌پذیری تست‌ها — را با پنج لایه‌ی معماری می‌سازد: لایه‌ی مکان‌یابی با نهانگاه دسته‌ها و جستجوی محدوده‌ای، لایه‌ی اشیای صفحه، لایه‌ی زبان مخصوص دامنه، کارخانه‌ی درایور با اتصال فازی، و لایه‌ی جداسازی داده و گزارش‌گیری. دو قرارداد پایداری (شناسه پایدار و قرارداد خوانش) رابطه‌ی برنامه‌ی هدف و چارچوب را قواعد‌مند می‌کند و تستِ نمودارهای گرافیکی را — که در درخت دسترس‌پذیری دیده نمی‌شوند — به تست داده‌ی خوانا تبدیل می‌کند.

ارزیابی روی برنامه‌ی نمونه‌ای پیچیده (جریان ورود چندپنجره‌ای، فیلترهای آبشاری، صفحه‌بندی، مجازی‌سازی جدول، نمودار و گزارش‌گیری) با ۲۵ سناریوی سرتاسری انجام شد؛ اجرای کامل پایدار، بازتولیدپذیر و در حدود شش دقیقه است و گزارش تعاملی
\lr{Allure}
با پیوست تصویر خطا تولید می‌شود. جداسازی پایگاه‌داده در سطح تست، استقلال و بازتولیدپذیری سناریوها را تضمین می‌کند و مسیر انتقال به
\lr{\gls{CI}}/\lr{\gls{CD}}
به‌صورت طراحی‌شده آماده است. همه‌ی اجزا متن‌باز و بدون هزینه‌ی مجوز هستند.

\section{کارهای آینده}
مسیر پیش‌روی چارچوب در چهار جهت است:
\begin{itemize}
	\item \textbf{اجرا در پایپلاین:} برپایی رانر ویندوزی، نهان‌گاه وابستگی‌ها و انتشار خودکار گزارش پس از هر تغییر؛ مسیر اجرای بی‌سر با
	\lr{DRIVER\_MODE=direct}
	آماده است.
	\item \textbf{گسترش پلتفرم:} افزودن درایورهای دیگر از الگوی کارخانه (
	\lr{Linux}
	با
	\lr{AT-SPI}،
	مک با
	\lr{Accessibility API})
	بدون تغییر لایه‌های بالا؛ قراردادهای پایداری مستقل از پلتفرم‌اند.
	\item \textbf{تولید خودکار تست:} استفاده از بازشماری کنترل‌های برنامه (ابزار بازرسی موجود در چارچوب) برای پیشنهاد اسکلت سناریوها و کشف انحراف شناسه‌ها در زمان ساخت برنامه.
	\item \textbf{تحلیل نتایج در گذر زمان:} نگهداری روند زمان اجرا و نرخ شکست اجراها برای آشکارسازی تست‌های شکننده پیش از آنکه تبدیل به مزاحم شوند.
\end{itemize}